%% ============================================================
%% STRESZCZENIE PO ANGIELSKU — OPIS FORMATKI
%% ============================================================
\addcontentsline{toc}{chapter}{Abstract}
\chapter*{Abstract}
This document is a \LaTeX{} template for theses prepared at the Faculty of Electronics, Telecommunications and Informatics, Gdańsk University of Technology. The structure of the thesis should comply with the guidelines specified in \href{\detokenize{https://cdn.files.pg.edu.pl/eti/Dziekanat/regulaminy/ZR%2045-2024.pdf}}{Rector’s Ordinance No.~45/2024} of 15 November 2024. These guidelines define the general rules for preparing theses and the editorial requirements. This template reflects, clarifies, and, in selected issues, extends these rules (including the numbering of tables, figures, equations, and citations). It also facilitates work with cross-references, captions, and the table of contents. It provides a preconfigured set of \LaTeX{} packages and practical guidance for the author, i.e., what to place in individual chapters, how to describe methods and results, and how to prepare figures and tables.

The main project file is \texttt{main.tex}. The formatting settings are collected in a separate class file \texttt{\_formatting.cls}, which contains all the required elements that build the thesis's structure. Additional formatting hints are described in separate chapters of the template. The document structure comprises four main chapters stored in the \emph{chapters} directory and an appendix with an example Python code. Each chapter is placed in a separate \texttt{.tex} file. The bibliography is managed with the \texttt{biblatex} (\texttt{biber}) package, in accordance with the ISO 690 standard, ensuring square-bracketed numbering in citation order. The title page and the student's declaration are attached as PDF files generated in the MyGUT (Moja PG) portal and inserted into the main file using the \verb|\includepdf| command.

The entire project was optimized for compilation in the Overleaf environment with the necessary use of the LuaLaTeX engine (TeX Live distribution, version at least 2022). The default font used in the form is Arial in size 10 pt, and its files are included in the \texttt{fonts} directory, which allows the project to be compiled in any environment. This approach combines automation of formatting with the flexibility of \LaTeX{}, ensuring compatibility with the standards of Gdańsk University of Technology. Please, remember that the document must meet the requirements of anti-plagiarism regulations. The generated PDF file cannot be graphic or encrypted in terms of text; it can also not use non-standard font encoding (e.g., by replacing character codes) or 'combining' diacritics from several characters into one glyph. Summary: this template aims to minimize the author's technical workload so they can focus on the content, while maintaining a uniform editorial style for all theses at the ETI Faculty. Any comments regarding the template should be addressed to its author, Ph.D. Eng. Sebastian Dziedziewicz, e-mail: \textit{sebastian.dziedziewicz@pg.edu.pl}. Below are example keywords that may be used in a thesis.

\section*{Keywords :}
\noindent machine learning, PID control, embedded systems, analog circuits, modulation, distributed systems
\section*{Field of science and technology in accordance with OECD requirements:}
\noindent engineering and technology; electrical, electronic and information engineering; robotics and automation; telecommunications.
